\section{Introduction}
\label{sec:introduction}
Recovery of coordinated electric vehicle (EV) charging requires knowing which
power limits can still execute. If a requested reduction is delayed, reallocating
its capacity elsewhere can overload the feeder. A later measurement of low draw
also need not exclude an unresolved increase or a higher fallback after profile
expiry. These ambiguities concern executable limits as well as communication
availability and charging demand.

For example, a 10-kW shared limit and active caps $(8,2)$ cannot support
reallocation to $(2,8)$ until the first reduction is established. If a temporary
2-kW profile later expires to an 8-kW fallback, the same risk returns even after
the reduction was observed. Recovery must account for both transitions.

Timed Command Interference links delayed commands and timeout fallback to EV
scheduling violations; its proposed future defenses include execution
verification and resilient coordination \cite{marandi2026tci}.
Related resilience mechanisms include capacity sharing
during communication outages \cite{gumrukcu2024dcs}, feedback scheduling under
uncertain charge acceptance \cite{kandpal2022feedback}, and distributed recovery
from corrupted scheduling messages \cite{qu2026resilient}. Networked predictive
control already addresses uncertain applied inputs and timestamped actuator
buffers \cite{pin2011networked}. These works establish the broader control and
communication context. Here the specific object is the feeder capacity reserved
for EV profiles that remain physically possible during recovery.

We construct a conservative ledger of active, unresolved, and expired profile
histories under an explicit trusted adapter contract. Before transmission, each
candidate command reserves its possible cap and subsequent fallback. Causal
post-application evidence can release predecessors while retaining later
unresolved generations. We compare two dispatch variants of this ledger:
per-group serialization (B3) and nonexpanding supersession (P).
Both share the same reservation accounting, scheduler, and information interface.
The contribution is this charging-specific ledger, its conditional affine
feasibility proposition, and a controlled evaluation of the variants' safety and
service tradeoffs by nonlinear AC replay. The support-function calculation and
generation fencing are established mechanisms. The adapter is a proposed
interface, not a guarantee supplied by an ordinary Open Charge Point Protocol
(OCPP) response
\cite{ocpp201ed4}.

\section{System and Execution Model}
\label{sec:model}
Each index $i$ denotes a feeder-connected charging phase group with actual
nonnegative power $p_i$. A common scheduler greedily assigns group caps in order
of earliest remaining vehicle departure, with group index breaking ties, subject
to the positive-part affine constraint below. Within a group, local
earliest-departure-first (EDF) allocation serves vehicles up to their constant acceptance
limits, remaining requests, and the effective group cap. Requests, arrivals,
departures, and acceptance limits are synthetic. An available booking service
reveals a request at arrival; delivered energy becomes known only through delayed
telemetry. This information interface is identical across policies and does not
reveal future arrivals or physical completion privately.

A command carries group identity, monotonically increasing generation, constant
cap $q_i$, and finite absolute validity interval $[s_i,e_i)$. The adapter has one
authenticated authority, enforces the effective cap, and maintains a persistent
known fallback $f_i$. Receiving a newer generation fences predecessor pending
operations; acceptance alone does not establish application. Application follows
the declared actuation lag and must finish before $e_i$. At absolute expiry, the
temporary profile switches to fallback without extending its lease. The model
uses exact clocks and zero tracking overshoot. Complete adapter state survives
modeled restart. This in-memory simulation establishes neither physical storage
durability nor charger implementation validity.

Channel impairments include denial, delay, and delayed application retries after
timeout or reconnection. They do not require reordered frames within an intact
TCP stream. Neither power telemetry nor command acceptance eliminates an
executable predecessor. A truthful application certificate identifies the group,
generation, effective cap, application time, absolute end, and fallback. Its
receipt is a separate delayed event.

Let $\mathcal H_i(t)$ conservatively retain cap histories consistent with initial
state, issued commands, elapsed time, and certificates received by $t$.
Let $c_i(h,s)$ denote the effective cap at time $s$ in history $h$, including
fallback, and let $r_i(I)$ be the power reserved for group $i$ throughout a future
interval $I\subseteq[t,\infty)$. The envelope and containment requirement are
\begin{equation}
 U_i(I)=\sup_{h\in\mathcal H_i(t),\,s\in I}c_i(h,s),
 \qquad 0\le p_i(s)\le U_i(I)\le r_i(I).
 \label{eq:envelope}
\end{equation}
The inequalities apply to every $s\in I$. The ledger can overapproximate
mutually incompatible histories. Its bound retains
an older high cap through uncertain reduction, every unresolved possible
increase, and fallback after each possible expiry, including fallback higher
than the temporary cap.

Finite differences about a declared operating point give signed affine rows
\begin{equation}
 A p+C d\le b,\qquad d\in[\underline d,\overline d],
 \label{eq:affine}
\end{equation}
where scalar $d$ scales background demand. Rows represent upper and lower
phase-voltage limits and source-rated transformer terminal-phase currents. The
source files provide no line ampacities, so line thermal feasibility is outside
this claim. Development margins tighten the affine limits; they are not certified
bounds on linearization error. With $X^+=\max(X,0)$ and $X^-=\min(X,0)$ entrywise,
the robust reservation test is
\begin{equation}
 A^+r+C^+\overline d+C^-\underline d\le b.
 \label{eq:robust}
\end{equation}
For each row, maximizing over the box $0\le p\le r$ assigns $p_i=r_i$ to
positive charging coefficients and zero to negative ones; maximizing the
background term selects the corresponding interval endpoint. Thus
\eqref{eq:robust} is the exact rowwise box counterpart of \eqref{eq:affine},
although the box can be conservative relative to correlated physical histories.

\section{Reservation Filter}
\label{sec:filter}
At each control event, the filter processes delivered evidence and advances its
ledger clock. It forms the nominal greedy targets from current bookings and
received energy estimates, then considers reductions before increases without
crediting the requested relief. For a candidate at group $i$, it constructs a
prospective envelope containing the existing histories, the candidate cap during
its lease, and fallback thereafter. All groups' current and candidate activation
and expiry breakpoints partition the full finite commitment span. Feasibility is
checked on every resulting interval, followed by the feasible persistent
fallback. An optimization or simulation horizon ending cannot release capacity.

For the supersession variant P, an envelope-expanding command is admitted only
when no unresolved operation
expands the evidence-only envelope of that group. Otherwise, a new command is
allowed only when its entire
validity and fallback span is nonexpanding relative to the existing envelope.
A smaller numerical cap alone is insufficient. This gate is recomputed after
each certificate because an older reduction certificate may reveal that a later
pending command is an expansion. The serialized variant B3 instead admits no
new command while that group has an unresolved operation. Both variants permit
different groups to progress concurrently.

The decision procedure is:
\begin{enumerate}
\setlength{\itemsep}{0pt}
\setlength{\parskip}{0pt}
\item Incorporate causal certificates and advance known absolute expiries.
\item Compute nominal targets and order candidate reductions before increases.
\item Partition all affected future commitments, including fallback transitions.
\item Limit the candidate using \eqref{eq:robust} on every interval and apply the
variant's unresolved-command gate.
\item Commit the prospective reservation before queueing the admitted command.
\item Query unresolved groups when available; release only histories excluded by
truthful evidence or verified expiry semantics.
\end{enumerate}

A certificate for generation $g$ removes predecessor possibilities through $g$,
but preserves all higher unresolved generations and their possible fallback
transitions. Expiry can remove an executable temporary command without removing
the fallback history it may have created. Reserve-before-send and retained state
are modeled logically; deployment requires an atomic durable implementation.

\begin{proposition}
\label{prop:reservation}
Assume initial containment, trusted cap enforcement and generation fencing,
causal truthful evidence, bounded background demand, and coverage of all future
commitments. If prospective reservations include each command and fallback,
satisfy \eqref{eq:robust}, and release histories only by the certificate and
expiry rules above, actual charging satisfies the selected affine rows
throughout every maintained feasible envelope.
\end{proposition}
\begin{proof}
Initially, the actual cap belongs to the retained history set. Reserving before
send includes both nondelivery and possible application. Local application,
fencing, and expiry follow covered histories; truthful evidence removes only
excluded histories and retains later possibilities. Induction over these events
therefore preserves $0\le p\le r$. The rowwise maximization yielding
\eqref{eq:robust} then bounds $Ap+Cd$ by $b$.
\end{proof}

This statement guarantees neither nonlinear AC feasibility nor minimum service
under indefinite outages, and does not establish recursive feasibility of future
service requests. AC outcomes are empirical. Proposition~\ref{prop:reservation}
applies to both ledger variants; neither is presumed to dominate in delivered
energy. In the primary experiment, both query uncertain state without
retransmitting lost commands.
